% Principe général de la génération augmentée par récupération
\begin{tikzpicture}[node distance=6mm and 11mm]

  \node[btpboxstrong, text width=2.5cm] (q) {Question\\utilisateur};
  \node[btpbox, text width=2.9cm, right=of q] (search)
    {Moteur de\\recherche\\[-1pt]\btplabel{lexical, sémantique, hybride}};
  \node[btpdb, right=of search, minimum width=2.3cm] (kb)
    {Base de\\connaissances};
  \node[btpbox, text width=2.9cm, below=16mm of search] (ctx)
    {Contexte\\[-1pt]\btplabel{extraits pertinents numérotés}};
  \node[btpbox, text width=2.5cm, left=of ctx] (prompt)
    {Invite\\[-1pt]\btplabel{question + contexte + consignes}};
  \node[btpboxalt, text width=2.9cm, right=of ctx] (llm)
    {Modèle de\\langage};
  \node[btpboxstrong, text width=2.7cm, below=14mm of llm] (rep)
    {Réponse\\avec citations};

  \draw[btparrow] (q) -- (search);
  \draw[btparrow] (search) -- (kb);
  \draw[btparrowlight] (kb.south) |- (ctx.east);
  \draw[btparrow] (ctx) -- (prompt);
  \draw[btparrow] (prompt.south) |- ($(rep.west)+(-0.7,0)$) -- (rep.west);
  \draw[btparrow] (ctx) -- (llm);
  \draw[btparrow] (llm) -- (rep);

  \node[btplabel, below=1mm of rep, text width=6cm]
    {Chaque affirmation est rattachable à un extrait source identifié};

\end{tikzpicture}
